\documentclass[11pt]{amsart}

\usepackage[margin=1.15in]{geometry}
\usepackage{amsmath,amssymb,amsthm,mathtools}
\usepackage{microtype}
\usepackage{hyperref}
\usepackage[nameinlink,capitalize]{cleveref}

\hypersetup{
  colorlinks=true,
  linkcolor=blue,
  citecolor=blue,
  urlcolor=blue
}

\theoremstyle{plain}
\newtheorem{theorem}{Theorem}[section]
\newtheorem{proposition}[theorem]{Proposition}
\newtheorem{lemma}[theorem]{Lemma}
\newtheorem{corollary}[theorem]{Corollary}

\theoremstyle{definition}
\newtheorem{definition}[theorem]{Definition}
\newtheorem{hypothesis}[theorem]{Hypothesis}

\theoremstyle{remark}
\newtheorem{remark}[theorem]{Remark}

\numberwithin{equation}{section}

\title[Compact ancient dynamics]{Compact Ancient Navier--Stokes Dynamics and Invariant Measures in Self-similar Variables}

\author{Guillem Duran Ballester}
\date{}

\subjclass[2020]{35Q30, 35B41, 35B44, 37A25, 76D05}
\keywords{Navier--Stokes equations, ancient solutions, invariant measures, self-similar variables, Type I singularities}

\begin{document}

\begin{abstract}
We study compact time-translation dynamics generated by bounded ancient
solutions of the three-dimensional Navier--Stokes equations in backward
self-similar variables.  The compactness used here is formulated in the
\(C^\infty_{\mathrm{loc}}\) topology: bounded ancient solutions generate
compact trajectory hulls, and uniform \(L^3\)-tightness is stable under
passage to such hulls.  On every compact invariant hull,
Krylov--Bogolyubov averages yield invariant Borel probability measures.  These
measures satisfy the stationary statistical weak form of the self-similar
Navier--Stokes equations.  Passing from this statistical identity to an exact
stationary equation for the barycenter is equivalent to the vanishing of an
explicit quadratic covariance.  Thus the invariant-measure construction is
separated from the additional rigidity needed to extract a single stationary
profile.
\end{abstract}

\maketitle

\section{Introduction}

We consider the incompressible Navier--Stokes equations in backward
self-similar variables,
\begin{equation}\label{eq:centered}
   \partial_\tau V+(V\cdot\nabla)V+\nabla P
   =
   \Delta V-\frac12(V+y\cdot\nabla V),
   \qquad \operatorname{div}V=0,
\end{equation}
on \(\mathbb R^3\times\mathbb R\).  The equation is autonomous in the
logarithmic time variable.  If \(V\) is a solution and \(s\in\mathbb R\), then
\[
   V^s(y,\tau)=V(y,\tau+s)
\]
is again a solution.  We study the compact dynamical system obtained by taking
the closure of these time translates in the \(C^\infty_{\mathrm{loc}}\)
topology.

The preceding minimal-element reduction identifies compact ancient dynamics as
the next object in the reduction: a nonzero normalized class, once closed under
local limits and time translation, contains a solution of least
\(L^\infty(\mathbb R^3\times\mathbb R)\) size.  The present paper does not
assert that such minimality alone proves recurrence or stationary behavior.
Instead it proves the compact dynamical consequences available once
one works on a compact trajectory hull.

Two issues must be kept distinct.  First, the invariant-measure construction
is a compactness argument on the hull: time averages of Dirac masses along an
orbit have weak limits, and every such limit is invariant under time
translation.  Second, the nonlinear term is retained exactly at the
statistical level: the stationary identity contains the average of
\(W\otimes W\), not an unspecified defect measure.  When this identity is
projected to the barycenter, an explicit quadratic covariance appears.
Vanishing of that covariance is an additional rigidity statement, and is the
precise bridge from invariant statistical dynamics to a single stationary
self-similar profile.

\subsection{Topological conventions}

All compactness statements are made in the Fr\'echet topology of local smooth
convergence, i.e. the \(C^\infty_{\mathrm{loc}}\) topology.  On compact
trajectory hulls this topology is metrizable, so sequential compactness and
compactness are equivalent.  The time-translation maps are homeomorphisms on
each invariant hull.  These elementary points are used below to pass from
orbit averages to invariant measures and from invariant measures to statements
about their supports.

For probability measures on a compact metric space we use the usual weak
topology, namely convergence against all continuous test functions.  The
compactness and sequential compactness of the space of Borel probability
measures in this topology are standard consequences of the Riesz
representation theorem and Banach--Alaoglu; see Billingsley
\cite{Billingsley1999}.  The time-average construction is the classical
Kryloff--Bogoliouboff argument \cite{KryloffBogoliouboff1937}.

\subsection{Results imported from the earlier papers}

We use the analytic reduction developed in the preceding papers of the series
without reproving it.  Paper~0 records the
pressure-normalized Caffarelli--Kohn--Nirenberg density, its critical scaling,
and the local epsilon-regularity mechanism used to detect nonzero blow-up
limits \cite[definition ``Pressure-normalized critical quantities,''
proposition ``Local finiteness and scaling,'' and theorem
``Caffarelli--Kohn--Nirenberg epsilon regularity'']{DuranPaper0}.
Paper~I constructs centered bounded ancient Type I profiles from Type I
blow-up, fixes the pressure gauges, proves nontriviality, and derives the
self-similar equation \eqref{eq:centered}
\cite[lemma ``Pressure gauges and local pressure bounds,'' proposition
``Ancient suitable Type I limit,'' proposition ``Nonzero ancient limit,'' and
lemma ``Renormalized equation'']{DuranPaper1}.

The present paper begins after those reductions.  Its PDE input is the bounded
mild ancient theory from Paper~III: the Ornstein--Uhlenbeck
semigroup formulation, the whole-space Helmholtz--Leray pressure convention,
local smoothing, smooth representatives, compactness of bounded families,
closure of the mild formulation under locally smooth convergence, invariance
under time translation, and closedness under uniform \(L^3\)-tightness
\cite[Section 2, definition ``Bounded mild ancient solution,'' lemma
``Local smoothing,'' remark ``smooth representatives,'' proposition ``Local
compactness,'' Section 3, and proposition ``Closedness under uniform
tightness'']{DuranPaper3}.  Every use below of local regularity, pressure
reconstruction, compactness of bounded families, or \(L^3\)-tightness is a
direct invocation of those results.

\section{Ancient solutions and trajectory hulls}

The terminology in this section is that of Paper~III
\cite[Section 2]{DuranPaper3}.  After applying the Helmholtz--Leray
projection \(\mathbb P\), \eqref{eq:centered} is interpreted as
\[
   \partial_\tau V
   =
   \Delta V-\frac12y\cdot\nabla V-\frac12V
   -\mathbb P\operatorname{div}(V\otimes V),
\]
and the mild formulation is written with the Ornstein--Uhlenbeck semigroup
generated by \(\Delta-\frac12y\cdot\nabla-\frac12\).  The semigroup formula,
the adjoint \(L^1\) bounds, and the differentiated kernel estimates used to
pass to locally smooth limits are those proved in
\cite[Section 2, Ornstein--Uhlenbeck semigroup and kernel estimates]{DuranPaper3}.
We write this semigroup as \(S(a)\), \(a>0\).

\begin{definition}[Bounded mild ancient solution]\label{def:mild}
A vector field \(V:\mathbb R^3\times\mathbb R\to\mathbb R^3\) is a bounded mild
ancient solution of \eqref{eq:centered} if
\[
   V\in L^\infty(\mathbb R^3\times\mathbb R),\qquad \operatorname{div}V=0
\]
in distributions, and if for every \(\sigma<\tau\)
\begin{equation}\label{eq:mild}
   V(\tau)
   =
   S(\tau-\sigma)V(\sigma)
   -
   \int_\sigma^\tau
   S(\tau-s)\mathbb P\operatorname{div}(V\otimes V)(s)\,ds
\end{equation}
as an identity in the weak-* topology of \(L^\infty(\mathbb R^3)\).
\end{definition}

This is the definition used in Paper~III
\cite[definition ``Bounded mild ancient solution'']{DuranPaper3}; it is
repeated only to fix notation for the dynamical hulls below.

\begin{definition}[Locally smooth convergence]\label{def:local-smooth}
A sequence \(W_n\) converges locally smoothly to \(W\) on
\(\mathbb R^3\times\mathbb R\) if, for every \(R,T<\infty\) and every integer
\(m\ge0\),
\[
   W_n\to W
   \quad\hbox{in }C^m(B_R\times[-T,T]).
\]
We use this locally smooth topology throughout.  It is the usual Fr\'echet
topology on \(C^\infty_{\mathrm{loc}}(\mathbb R^3\times\mathbb R)\), and in
particular it is metrizable.
\end{definition}

Explicitly, one may choose a countable exhaustion by compact cylinders
\[
   Q_N=\overline{B_N}\times[-N,N],\qquad N\ge1,
\]
and use the seminorms
\[
   p_{N,m}(W)
   =
   \sum_{|\alpha|+j\le m}
   \sup_{(y,\tau)\in Q_N}
   |\partial_y^\alpha\partial_\tau^j W(y,\tau)|.
\]
The metric
\[
   d(W,Z)=
   \sum_{N,m\ge1}2^{-N-m}
   \frac{p_{N,m}(W-Z)}{1+p_{N,m}(W-Z)}
\]
generates the same topology.  Therefore, on any compact set \(K\) in this
topology, every seminorm \(p_{N,m}\) is uniformly bounded on \(K\), and the
evaluation maps
\[
   W\mapsto \partial_y^\alpha\partial_\tau^jW(y,\tau)
\]
are continuous for fixed \((y,\tau)\).  If \(B_R\times[-T,T]\) is fixed, then
the restriction map
\[
   W\mapsto W|_{B_R\times[-T,T]}
\]
is continuous from \(K\) into \(C^m(B_R\times[-T,T])\).  Hence every
integral over a fixed compact spatial set whose integrand is a polynomial in
finitely many derivatives of \(W\), multiplied by a fixed smooth compactly
supported test function, defines a continuous real-valued function on \(K\).
These elementary facts are the only topological input used in the sequel.

\begin{lemma}[Local smoothing and pressure reconstruction]\label{lem:classical-local}
Let \(V\) be a bounded mild solution of \eqref{eq:centered} on
\(\mathbb R^3\times(\tau_0,\tau_1)\).  Then \(V\) has the smooth
representative supplied by Paper~III.  For every compact cylinder
\(Q\Subset\mathbb R^3\times(\tau_0,\tau_1)\) there is a smooth pressure
function \(P_Q\), unique up to addition of a function of time, such that
\((V,P_Q)\) satisfies \eqref{eq:centered} pointwise on \(Q\).
\end{lemma}

\begin{proof}
Local smoothing and uniform \(C^m\)-bounds for bounded mild solutions are
\cite[lemma ``Local smoothing'']{DuranPaper3}.  The passage from the projected
equation to the local pressure \(P_Q\), with the whole-space
Helmholtz--Leray pressure normalization and uniqueness up to an additive
function of time, is the pressure convention and smooth-representative
statement in \cite[Section 2 and remark ``Representatives and pointwise
bounds'']{DuranPaper3}.  No additional regularity or decay assumption is
introduced here.
\end{proof}

\begin{proposition}[Compactness of bounded families]\label{prop:bounded-compactness}
Let \(V_n\) be bounded mild ancient solutions of \eqref{eq:centered} satisfying
\[
   \sup_n\|V_n\|_{L^\infty(\mathbb R^3\times\mathbb R)}<\infty .
\]
Then a subsequence converges locally smoothly to a bounded mild ancient
solution of \eqref{eq:centered}.
\end{proposition}

\begin{proof}
This is the compactness and closedness theorem proved in
\cite[proposition ``Local compactness'']{DuranPaper3}: the local smoothing
estimates give Arzela--Ascoli compactness on compact cylinders, and the
semigroup/adjoint kernel estimates pass the mild formulation to the locally
smooth limit.
\end{proof}

\begin{theorem}[Compact trajectory hull]\label{thm:compact-hull}
Let \(V\) be a bounded mild ancient solution of \eqref{eq:centered}.  Let \(K\)
be the closure, in the locally smooth topology, of the family
\[
   \{V^s:\ s\in\mathbb R\}.
\]
Then \(K\) is compact.  Every \(W\in K\) is a bounded mild ancient solution of
\eqref{eq:centered}, with
\[
   \|W\|_{L^\infty(\mathbb R^3\times\mathbb R)}
   \le
   \|V\|_{L^\infty(\mathbb R^3\times\mathbb R)}.
\]
Moreover the maps
\[
   \Theta_t W=W^t,\qquad t\in\mathbb R,
\]
define a jointly continuous group of transformations of \(K\).

If, in addition, \(V\) satisfies
\begin{equation}\label{eq:L3-tight}
   \sup_{\tau\in\mathbb R}\|V(\tau)\|_{L^3(\mathbb R^3)}\le L
\end{equation}
and there is \(\omega:[1,\infty)\to[0,\infty)\), with
\(\omega(R)\to0\) as \(R\to\infty\), such that
\begin{equation}\label{eq:tail}
   \sup_{\tau\in\mathbb R}
   \int_{|y|>R}|V(y,\tau)|^3\,dy\le\omega(R)
   \qquad R\ge1,
\end{equation}
then every \(W\in K\) satisfies the corresponding estimates
\[
   \sup_{\tau\in\mathbb R}\|W(\tau)\|_{L^3(\mathbb R^3)}\le L
\]
and
\[
   \sup_{\tau\in\mathbb R}
   \int_{|y|>R}|W(y,\tau)|^3\,dy\le\omega(R),
   \qquad R\ge1 .
\]
\end{theorem}

\begin{proof}
Every translate \(V^s\) satisfies the same mild equation and the same
\(L^\infty\) bound as \(V\), because \eqref{eq:centered} and
\eqref{eq:mild} are autonomous in \(\tau\).  More explicitly, replacing
\(\sigma,\tau\) in \eqref{eq:mild} by \(\sigma+s,\tau+s\) gives the mild
identity for \(V^s\), and the essential supremum is unchanged by the time
translation.  Therefore every sequence \(V^{s_n}\) has, by
\cref{prop:bounded-compactness}, a locally smoothly convergent subsequence
whose limit is a bounded mild ancient solution.

We now prove compactness of the closure \(K\) directly.  Let \((W_n)\) be an
arbitrary sequence in \(K\).  Since \(K\) is the closure of the orbit, for
each \(n\) there exists \(s_n\in\mathbb R\) such that
\[
   d(W_n,V^{s_n})<\frac1n,
\]
where \(d\) is the metric defined above for the locally smooth topology.  The
sequence \((V^{s_n})\) has a locally smoothly convergent subsequence
\((V^{s_{n_k}})\), say
\[
   V^{s_{n_k}}\to W
   \qquad\hbox{in }C^\infty_{\mathrm{loc}}(\mathbb R^3\times\mathbb R).
\]
Since \(d(W_{n_k},V^{s_{n_k}})\le n_k^{-1}\to0\), the triangle inequality
gives \(W_{n_k}\to W\) in the same topology.  The set \(K\) is closed by
definition, hence \(W\in K\).  Thus every sequence in \(K\) has a convergent
subsequence with limit in \(K\).  Because the locally smooth topology is
metrizable, sequential compactness is equivalent to compactness, and \(K\) is
compact.

If \(W\in K\), then there are \(s_n\in\mathbb R\) with \(V^{s_n}\to W\)
locally smoothly.  For every point \((y,\tau)\),
\[
   |W(y,\tau)|
   =
   \lim_{n\to\infty}|V(y,\tau+s_n)|
   \le
   \|V\|_{L^\infty(\mathbb R^3\times\mathbb R)}.
\]
Here the \(L^\infty\) norm is the essential supremum, but the smooth
representatives supplied by \cref{lem:classical-local} obey the same bound
pointwise: if a continuous representative exceeded the essential supremum at
one point, it would exceed it on a nonempty open set.  Thus the displayed
\(L^\infty\) estimate holds.  The proof of
\cref{prop:bounded-compactness}, equivalently
\cite[proposition ``Local compactness'']{DuranPaper3}, applied to this
convergent sequence of translates, also shows that \(W\) satisfies the mild
formulation and is divergence free.

The identity \(\Theta_t\Theta_s=\Theta_{t+s}\) follows from the definition of
time translation.  The map \(\Theta_t\) sends the orbit to itself; the
continuity proved below then implies that it sends its closure \(K\) into
\(K\).  If \(W_n\to W\) locally smoothly, then for every compact
cylinder \(B_R\times[-T,T]\) and every \(m\),
\[
   \|\Theta_tW_n-\Theta_tW\|_{C^m(B_R\times[-T,T])}
   =
   \|W_n-W\|_{C^m(B_R\times[-T+t,T+t])},
\]
where the right-hand side is estimated on \(B_R\times[-T-|t|,T+|t|]\), one of
the compact cylinders used in the locally smooth topology.  The right-hand
side tends to zero, so \(\Theta_t\) is
continuous.  Applying the same argument to \(\Theta_{-t}\) gives
\(\Theta_{-t}\Theta_t=\Theta_t\Theta_{-t}=\operatorname{Id}\) on \(K\).  Hence
each \(\Theta_t\) is a homeomorphism of \(K\).

It remains to verify joint continuity of the action.  Suppose \(t_n\to t\) and
\(W_n\to W\) locally smoothly.  Fix \(R,T,m\).  Choose \(A>|t|+1\) such that
\(|t_n|\le A\) for all large \(n\).  Then
\begin{align*}
   \|\Theta_{t_n}W_n-\Theta_tW\|_{C^m(B_R\times[-T,T])}
   &\le
   \|\Theta_{t_n}W_n-\Theta_{t_n}W\|_{C^m(B_R\times[-T,T])}
   \\
   &\quad+
   \|\Theta_{t_n}W-\Theta_tW\|_{C^m(B_R\times[-T,T])}.
\end{align*}
The first term is controlled by the \(C^m\)-norm of \(W_n-W\) on
\(B_R\times[-T-A,T+A]\), hence tends to zero.  The second term tends to zero
because \(W\) is smooth, so each derivative
\(\partial_y^\alpha\partial_\tau^jW\), \(|\alpha|+j\le m\), is uniformly
continuous on the compact cylinder \(B_R\times[-T-A,T+A]\).  Thus
\((t,W)\mapsto\Theta_tW\) is continuous, and the maps form a continuous group.

It remains to prove the uniform \(L^3\) assertions.  The family of translates
\(\{V^s:s\in\mathbb R\}\) satisfies \eqref{eq:L3-tight} and \eqref{eq:tail}
with the same constants because these estimates are invariant under
\(\tau\)-translation.  Since \(K\) is the locally smooth closure of this
family, the closedness under uniform \(L^3\)-tightness from
\cite[proposition ``Closedness under uniform tightness'']{DuranPaper3}
applies directly.  It gives the two displayed \(L^3\) estimates for every
\(W\in K\), with the same \(L\) and the same tail modulus \(\omega\).  No
additional tightness, decay, or compact-support hypothesis is introduced.
\end{proof}

\section{Compact invariant sets}

\begin{definition}[Compact invariant set]\label{def:compact-invariant}
A nonempty compact subset \(K\), with respect to the locally smooth topology,
of bounded mild ancient solutions is invariant if \(\Theta_tK=K\) for every
\(t\in\mathbb R\).  It is minimal if it has no proper nonempty compact
invariant subset.
\end{definition}

\begin{lemma}[Continuity of the restricted translation action]\label{lem:action-continuity}
Let \(K\) be a compact invariant set in the sense of
\cref{def:compact-invariant}.  Then, for every \(t\in\mathbb R\),
\(\Theta_t:K\to K\) is a homeomorphism, and the action
\[
   \mathbb R\times K\ni(t,W)\mapsto\Theta_tW\in K
\]
is jointly continuous.
\end{lemma}

\begin{proof}
The identity \(\Theta_t\Theta_s=\Theta_{t+s}\) follows from the definition of
time translation.  Since \(K\) is invariant, \(\Theta_tK=K\) for every
\(t\), so \(\Theta_t\) maps \(K\) onto \(K\) and has inverse
\(\Theta_{-t}\).

Fix \(t\in\mathbb R\).  If \(W_n\to W\) in \(K\) locally smoothly, then for
every \(R,T<\infty\) and every integer \(m\ge0\),
\[
   \|\Theta_tW_n-\Theta_tW\|_{C^m(B_R\times[-T,T])}
   =
   \|W_n-W\|_{C^m(B_R\times[-T+t,T+t])}.
\]
The interval \([-T+t,T+t]\) is contained in \([-T-|t|,T+|t|]\), so the
right-hand side tends to zero by local smooth convergence.  Thus
\(\Theta_t\) is continuous on \(K\).  Applying the same argument to
\(\Theta_{-t}\) proves that \(\Theta_t\) is a homeomorphism.

It remains to prove joint continuity.  Suppose \(t_n\to t\) and
\(W_n\to W\) in \(K\).  Fix \(R,T,m\), and choose \(A>|t|+1\) so that
\(|t_n|\le A\) for all sufficiently large \(n\).  Then
\begin{align*}
   \|\Theta_{t_n}W_n-\Theta_tW\|_{C^m(B_R\times[-T,T])}
   &\le
   \|\Theta_{t_n}W_n-\Theta_{t_n}W\|_{C^m(B_R\times[-T,T])} \\
   &\quad+
   \|\Theta_{t_n}W-\Theta_tW\|_{C^m(B_R\times[-T,T])}.
\end{align*}
The first term is bounded by the \(C^m\)-norm of \(W_n-W\) on
\(B_R\times[-T-A,T+A]\), and hence tends to zero.  For the second term,
each derivative \(\partial_y^\alpha\partial_\tau^jW\), with
\(|\alpha|+j\le m\), is uniformly continuous on
\(\overline B_R\times[-T-A,T+A]\).  Therefore
\[
   \partial_y^\alpha\partial_\tau^jW(y,\tau+t_n)
   \to
   \partial_y^\alpha\partial_\tau^jW(y,\tau+t)
\]
uniformly for \((y,\tau)\in\overline B_R\times[-T,T]\).  The second term
also tends to zero, proving joint continuity.
\end{proof}

\begin{proposition}[Minimal compact subsets]\label{prop:minimal-subset}
Every nonempty compact invariant set contains a minimal compact invariant
subset.
\end{proposition}

\begin{proof}
Let \(\mathcal F\) be the family of nonempty compact invariant subsets of
\(K\).  We order \(\mathcal F\) by reverse inclusion:
\[
   A\preceq B
   \quad\Longleftrightarrow\quad
   B\subset A .
\]
A maximal element for this order is exactly a minimal element for ordinary
set inclusion.  To apply Zorn's lemma, let \(\{K_\alpha\}\) be a
\(\preceq\)-chain.  Equivalently, it is a family totally ordered by ordinary
inclusion.  The intersection
\[
   K_*:=\bigcap_\alpha K_\alpha
\]
is nonempty: each \(K_\alpha\) is closed in the ambient compact set \(K\), and
finite intersections are nonempty because a finite subfamily of a totally
ordered family by inclusion has a smallest member.  Thus the finite
intersection property and compactness of \(K\) give \(K_*\ne\emptyset\).
The set \(K_*\) is compact because it is closed in \(K\).

It is invariant by \cref{lem:action-continuity}.  Indeed, since every member
of the chain is invariant,
\[
   \Theta_tK_*
   \subset
   \bigcap_\alpha \Theta_tK_\alpha
   =
   K_* .
\]
Applying the same inclusion with \(-t\) gives
\[
   \Theta_{-t}K_*
   \subset
   K_* .
\]
Applying \(\Theta_t\) to this last inclusion and using
\(\Theta_t\Theta_{-t}=\operatorname{Id}\) gives the reverse inclusion
\[
   K_*
   \subset
   \Theta_tK_* .
\]
Hence \(\Theta_tK_*=K_*\).  Therefore \(K_*\in\mathcal F\).  Moreover
\(K_*\subset K_\alpha\) for every \(\alpha\), so \(K_\alpha\preceq K_*\) for
every \(\alpha\); that is, \(K_*\) is an upper bound for the chain in the
reverse-inclusion order.  Zorn's lemma gives a maximal element of
\((\mathcal F,\preceq)\), equivalently a nonempty compact invariant subset
which is minimal under ordinary inclusion.
\end{proof}

\begin{lemma}[The zero trajectory]\label{lem:zero-trajectory}
Let \(K\) be a minimal compact invariant set.  If the zero solution belongs to
\(K\), then \(K\) consists only of the zero solution.
\end{lemma}

\begin{proof}
The zero solution satisfies \(\Theta_t0=0\) for every \(t\), so the singleton
\(\{0\}\) is compact and invariant.  If \(0\in K\), then \(\{0\}\) is a
nonempty compact invariant subset of \(K\).  Minimality of \(K\) therefore
forces \(K=\{0\}\).
\end{proof}

\section{Invariant measures}

\begin{theorem}[Krylov--Bogolyubov averages]\label{thm:KB}
Let \(K\) be a compact invariant set and let \(W_0\in K\).  For \(T>0\), define
the probability measure
\[
   \mu_T=\frac1T\int_0^T\delta_{\Theta_sW_0}\,ds
\]
on \(K\).  If \(T_n\to\infty\), then a subsequence of \(\mu_{T_n}\) converges
weakly to a Borel probability measure \(\mu\) on \(K\).  Every such limit is
invariant:
\[
   \int_K F(\Theta_tW)\,d\mu(W)=\int_K F(W)\,d\mu(W)
\]
for every \(t\in\mathbb R\) and every continuous function \(F:K\to\mathbb R\).
If \(K\) is minimal, then the support of every invariant Borel probability
measure on \(K\) is equal to \(K\).
\end{theorem}

\begin{proof}
By \cref{lem:action-continuity}, the orbit map \(s\mapsto\Theta_sW_0\) is
continuous.  The formula defining \(\mu_T\) denotes the push-forward of
normalized Lebesgue measure on \([0,T]\) under this orbit map,
\[
   [0,T]\ni s\mapsto \Theta_sW_0\in K .
\]
Thus \(\mu_T\) is a Borel probability measure on \(K\).  In particular, for a
Borel set \(A\subset K\),
\[
   \mu_T(A)=\frac1T\left|\{s\in[0,T]:\Theta_sW_0\in A\}\right|,
\]
where the set on the right is Lebesgue measurable because the orbit map is
continuous and \(A\) is Borel.

For \(F\in C(K)\), the map
\[
   s\mapsto F(\Theta_sW_0)
\]
is continuous, hence Borel and bounded on every finite interval.  Therefore
\(\mu_T\) is the Borel probability measure characterized by
\[
   \int_K F(W)\,d\mu_T(W)
   =
   \frac1T\int_0^TF(\Theta_sW_0)\,ds .
\]
Since \(K\) is compact and metrizable, \(C(K)\) is separable.  By the Riesz
representation theorem, Borel probability measures on \(K\) are the positive
norm-one elements of \(C(K)^*\).  Banach--Alaoglu gives weak-* compactness of
the unit ball of \(C(K)^*\), and separability of \(C(K)\) makes this compact
set metrizable on bounded subsets.  Consequently every sequence of Borel
probability measures on \(K\) has a weakly convergent subsequence
\cite{Billingsley1999}.  Applying this to \(\mu_{T_n}\), choose a subsequence
which converges weak-* in \(C(K)^*\) to a bounded linear functional \(L\) on
\(C(K)\).  In concrete terms,
\[
   \int_K F\,d\mu_{T_{n_k}}\to L(F)
   \qquad\hbox{for every }F\in C(K).
\]
For every \(F\in C(K)\) with \(F\ge0\), one has
\[
   L(F)=\lim_{k\to\infty}\int_K F\,d\mu_{T_{n_k}}\ge0,
\]
and, for the constant function \(1\),
\[
   L(1)=\lim_{k\to\infty}\mu_{T_{n_k}}(K)=1 .
\]
The Riesz representation theorem therefore identifies \(L\) with integration
against a Borel probability measure \(\mu\) on \(K\).  Thus the subsequence
converges weakly to \(\mu\).

Let \(F\) be continuous and fix \(t>0\).  Then
\[
   \int_K F(\Theta_tW)\,d\mu_T(W)-\int_K F(W)\,d\mu_T(W)
   =
   \frac1T\int_T^{T+t}F(\Theta_sW_0)\,ds
   -
   \frac1T\int_0^tF(\Theta_sW_0)\,ds .
\]
The absolute value of the right-hand side is at most
\[
   \frac{2t}{T}\|F\|_{L^\infty(K)}.
\]
The displayed equality follows by the change of variables \(r=s+t\) in the
first time average.  Since \(W\mapsto F(\Theta_tW)\) is continuous on \(K\),
weak convergence of \(\mu_{T_{n_k}}\) to \(\mu\) gives convergence of both
integrals on the left along the chosen subsequence.  Letting
\(T=T_{n_k}\to\infty\) gives invariance for \(t>0\).  If \(t<0\), apply the
positive-time result with \(-t\) to the continuous function
\(F\circ\Theta_t\).  Since
\(\Theta_{-t}\Theta_t\) is the identity, this gives the asserted identity for
negative \(t\) as well.  For \(t=0\) the identity is trivial.  Equivalently,
\((\Theta_t)_\#\mu=\mu\) as Borel measures for every \(t\in\mathbb R\).  The
last equivalence follows from the uniqueness part of the Riesz representation
theorem, because the two probability measures have the same integral against
every \(F\in C(K)\).

If \(K\) is minimal and \(\mu\) is invariant, let
\(\operatorname{supp}\mu\) denote the closed support, namely the set of all
points \(W\in K\) such that every open neighborhood of \(W\) has positive
\(\mu\)-measure.  Its complement is open: if
\(W\notin\operatorname{supp}\mu\), then some open neighborhood \(U\) of \(W\)
has \(\mu(U)=0\), and the same \(U\) is a zero-measure neighborhood for each
of its points.  Hence \(\operatorname{supp}\mu\) is closed.  This set is
nonempty: if it were empty, then every \(W\in K\) would have an open
neighborhood \(U_W\) with \(\mu(U_W)=0\); compactness would give a finite
subcover of \(K\), and hence \(\mu(K)=0\), contradicting \(\mu(K)=1\).  The
support is compact because it is closed in the compact space \(K\).

It is also invariant.  Indeed, if \(W\in\operatorname{supp}\mu\) and \(U\) is
an open neighborhood of \(\Theta_tW\), then \(\Theta_{-t}U\) is an open
neighborhood of \(W\), hence \(\mu(\Theta_{-t}U)>0\).  Since
\((\Theta_t)_\#\mu=\mu\),
\[
   \mu(U)=\mu(\Theta_{-t}U)>0 .
\]
It follows that \(\Theta_tW\in\operatorname{supp}\mu\), so
\(\Theta_t(\operatorname{supp}\mu)\subset\operatorname{supp}\mu\).  Applying
the same argument to \(-t\) gives the opposite inclusion.  Hence the support
is invariant.  The support is therefore a nonempty compact invariant subset
of \(K\), and minimality forces it to be \(K\).
\end{proof}

\begin{corollary}[Nonzero support]\label{cor:nonzero-support}
Let \(K\) be a minimal compact invariant set containing at least one nonzero
trajectory.  Then the zero trajectory is not in \(K\), and every invariant
Borel probability measure on \(K\) has support equal to \(K\).  In particular,
no such measure is supported at the zero solution.
\end{corollary}

\begin{proof}
The first assertion follows from \cref{lem:zero-trajectory}.  The support
statement is the last assertion of \cref{thm:KB}.
\end{proof}

\section{The stationary statistical identity}

The next theorem records the exact weak identity satisfied by invariant
measures.  The test fields are solenoidal, so the pressure is eliminated in
the usual distributional formulation.

\begin{theorem}[Stationary statistical form]\label{thm:statistical-identity}
Let \(K\) be a compact invariant set of bounded mild ancient solutions of
\eqref{eq:centered}, and let \(\mu\) be an invariant Borel probability measure
on \(K\).  Then, for every solenoidal test field
\(\phi\in C_c^\infty(\mathbb R^3;\mathbb R^3)\),
\begin{equation}\label{eq:statistical-identity}
   \begin{aligned}
   0
   =
   \int_K
   \bigg[
      &\int_{\mathbb R^3}
      W(y,0)\cdot
      \left(\Delta\phi(y)+\phi(y)+\frac12 y\cdot\nabla\phi(y)\right)\,dy \\
      &+
      \int_{\mathbb R^3}
      W_i(y,0)W_j(y,0)\partial_j\phi_i(y)\,dy
   \bigg]d\mu(W).
   \end{aligned}
\end{equation}
\end{theorem}

\begin{proof}
For such a test field define
\[
   F(W)=\int_{\mathbb R^3}W(y,0)\cdot\phi(y)\,dy .
\]
If \(W_n\to W\) locally smoothly, then \(W_n(\cdot,0)\to W(\cdot,0)\)
uniformly on the compact support of \(\phi\).  Hence \(F(W_n)\to F(W)\), so
\(F\in C(K)\).
Invariance gives
\[
   \int_K F(\Theta_hW)\,d\mu(W)=\int_K F(W)\,d\mu(W)
   \qquad h\in\mathbb R .
\]
Therefore, for \(0<|h|\le1\),
\[
   \int_K \frac{F(\Theta_hW)-F(W)}{h}\,d\mu(W)=0 .
\]
Choose \(R\) such that \(\operatorname{supp}\phi\subset B_R\).  Because
\[
   W\mapsto \|\partial_\tau W\|_{L^\infty(\overline B_R\times[-1,1])}
\]
is continuous on \(K\) and \(K\) is compact, the quantity
\[
   M_\phi:=
   \sup_{W\in K}\|\partial_\tau W\|_{L^\infty(\overline B_R\times[-1,1])}
\]
is finite.  For \(0<|h|\le1\), the fundamental theorem of calculus in the
\(\tau\)-variable gives, pointwise for \(y\in\operatorname{supp}\phi\),
\[
   \frac{F(\Theta_hW)-F(W)}{h}
   =
   \int_{\mathbb R^3}
   \left(\frac1h\int_0^h\partial_\tau W(y,s)\,ds\right)\cdot\phi(y)\,dy .
\]
The absolute value is bounded by
\[
   |\operatorname{supp}\phi|\,\|\phi\|_{L^\infty} M_\phi,
\]
which is independent of \(W\) and \(h\).  For each fixed \(W\), smoothness in
\(\tau\) gives convergence of the inner average to
\(\partial_\tau W(y,0)\), uniformly for \(y\in\operatorname{supp}\phi\).
The integrands as functions of \(W\) are Borel, because they are continuous
with respect to the locally smooth topology.  Dominated convergence on
\((K,\mu)\) gives
\[
   \int_K\int_{\mathbb R^3}\partial_\tau W(y,0)\cdot\phi(y)\,dy\,d\mu(W)=0 .
\]
For each \(W\in K\), \cref{lem:classical-local} gives a smooth local pressure
on a cylinder containing \(\operatorname{supp}\phi\times\{0\}\).  We write it
as \(P\) on that cylinder.  Pairing \eqref{eq:centered} at time \(0\) with
\(\phi\), in the sense of distributions and then classically on the compact
support of \(\phi\), gives
\[
   \begin{aligned}
   \int_{\mathbb R^3}\partial_\tau W\cdot\phi\,dy
   =
   &\int_{\mathbb R^3}\Delta W\cdot\phi\,dy
   -\frac12\int_{\mathbb R^3}W\cdot\phi\,dy
   -\frac12\int_{\mathbb R^3}(y\cdot\nabla W)\cdot\phi\,dy \\
   &-\int_{\mathbb R^3}(W\cdot\nabla)W\cdot\phi\,dy
   -\int_{\mathbb R^3}\nabla P\cdot\phi\,dy .
   \end{aligned}
\]
All terms in this display are evaluated at \(\tau=0\), and all integrations by parts below
are over compact subsets of the pressure cylinder, because \(\phi\) is
compactly supported.  The pressure term vanishes:
\[
   \int_{\mathbb R^3}\nabla P\cdot\phi\,dy
   =
   -\int_{\mathbb R^3}P\,\operatorname{div}\phi\,dy=0 .
\]
Two integrations by parts give
\[
   \int_{\mathbb R^3}\Delta W\cdot\phi\,dy
   =
   \int_{\mathbb R^3}W\cdot\Delta\phi\,dy .
\]
Using \(\operatorname{div}W=0\), the convection term is
\[
   -\int_{\mathbb R^3}(W\cdot\nabla)W\cdot\phi\,dy
   =
   -\int_{\mathbb R^3}W_j\partial_jW_i\,\phi_i\,dy
   =
   \int_{\mathbb R^3}W_iW_j\partial_j\phi_i\,dy .
\]
Here and below repeated spatial indices are summed over \(\{1,2,3\}\).
For the drift term,
\[
   \int_{\mathbb R^3}(y\cdot\nabla W)\cdot\phi\,dy
   =
   \int_{\mathbb R^3}y_j\partial_jW_i\,\phi_i\,dy
   =
   -\int_{\mathbb R^3}W_i\partial_j(y_j\phi_i)\,dy .
\]
Since \(\partial_jy_j=3\), this equals
\[
   -3\int_{\mathbb R^3}W\cdot\phi\,dy
   -
   \int_{\mathbb R^3}W\cdot(y\cdot\nabla\phi)\,dy .
\]
Therefore
\[
   -\frac12\int_{\mathbb R^3} W\cdot\phi\,dy
   -
   \frac12\int_{\mathbb R^3} (y\cdot\nabla W)\cdot\phi\,dy
   =
   \int_{\mathbb R^3} W\cdot\phi\,dy
   +
   \frac12\int_{\mathbb R^3} W\cdot(y\cdot\nabla\phi)\,dy .
\]
Combining these identities, and evaluating \(W\) at \(\tau=0\), gives for
each \(W\)
\[
   \begin{aligned}
   \int_{\mathbb R^3}\partial_\tau W(y,0)\cdot\phi(y)\,dy
   =
   &\int_{\mathbb R^3}
   W(y,0)\cdot
   \left(\Delta\phi(y)+\phi(y)+\frac12 y\cdot\nabla\phi(y)\right)\,dy\\
   &+
   \int_{\mathbb R^3}
   W_i(y,0)W_j(y,0)\partial_j\phi_i(y)\,dy .
   \end{aligned}
\]
The right-hand side is a continuous function of \(W\in K\), by local smooth
convergence on \(\operatorname{supp}\phi\times\{0\}\).  It is bounded on
\(K\) because the seminorm
\[
   W\mapsto \|W(\cdot,0)\|_{L^\infty(\operatorname{supp}\phi)}
\]
is bounded on the compact set \(K\), and the displayed expression is bounded
by a constant depending only on \(\phi\), \(|\operatorname{supp}\phi|\), and
this uniform local \(L^\infty\) bound.  Therefore it is \(\mu\)-integrable.
Substituting the pointwise identity for
\(\int \partial_\tau W(y,0)\cdot\phi(y)\,dy\) into the averaged identity
\[
   \int_K\int_{\mathbb R^3}\partial_\tau W(y,0)\cdot\phi(y)\,dy\,d\mu(W)=0
\]
gives \eqref{eq:statistical-identity}.
\end{proof}

\begin{theorem}[Barycenter and covariance]\label{thm:barycenter}
Let \(K\) and \(\mu\) be as in \cref{thm:statistical-identity}.  Define the
barycenter
\[
   \overline W(y)=\int_K W(y,0)\,d\mu(W)
\]
and the quadratic covariance
\[
   C_{ij}(y)=
   \int_K W_i(y,0)W_j(y,0)\,d\mu(W)
   -
   \overline W_i(y)\overline W_j(y).
\]
Then \(\overline W\) and \(C\) are \(C^\infty_{\mathrm{loc}}\), and
\(\overline W\) is solenoidal.  Moreover, for every solenoidal
\(\phi\in C_c^\infty(\mathbb R^3;\mathbb R^3)\),
\begin{equation}\label{eq:barycenter-equation}
   \int_{\mathbb R^3}
   \overline W\cdot
   \left(\Delta\phi+\phi+\frac12 y\cdot\nabla\phi\right)\,dy
   +
   \int_{\mathbb R^3}
   \overline W_i\overline W_j\partial_j\phi_i\,dy
   +
   \int_{\mathbb R^3}
   C_{ij}\partial_j\phi_i\,dy
   =0 .
\end{equation}
Consequently, \(\overline W\) is a weak stationary solution of
\eqref{eq:centered}, in the pressure-free formulation against solenoidal
test fields, if and only if
\begin{equation}\label{eq:zero-covariance}
   \int_{\mathbb R^3}C_{ij}(y)\partial_j\phi_i(y)\,dy=0
\end{equation}
for every solenoidal
\(\phi\in C_c^\infty(\mathbb R^3;\mathbb R^3)\).
\end{theorem}

\begin{proof}
Fix a compact ball \(B_R\) and a multiindex \(\alpha\).  Compactness of \(K\)
in the locally smooth topology gives
\[
   M_{R,\alpha}
   :=
   \sup_{W\in K}\sup_{y\in B_R}
   |\partial_y^\alpha W(y,0)|
   <\infty .
\]
For each \(y\), the evaluation map \(W\mapsto W(y,0)\) is continuous from
\(K\) into the finite-dimensional space \(\mathbb R^3\).  Its image is compact,
hence bounded and Borel measurable.  The barycenter is therefore well defined
pointwise as a finite vector-valued Bochner integral; equivalently, since the
range is \(\mathbb R^3\), it is obtained by integrating the three scalar
components separately.  We next justify the spatial differentiation under the
integral sign.

Let \(e_\ell\) be a coordinate vector and let \(y\in B_R\).  If
\(|h|\) is small enough that the segment
\(\{y+\theta h e_\ell:0\le\theta\le1\}\) is contained in a fixed slightly
larger ball \(B_{R'}\), then
\[
   \frac{\partial_y^\alpha W(y+he_\ell,0)-\partial_y^\alpha W(y,0)}{h}
   =
   \int_0^1
   \partial_\ell\partial_y^\alpha W(y+\theta h e_\ell,0)\,d\theta .
\]
The right-hand side is bounded uniformly in \(W\in K\) by
\(M_{R',\alpha+e_\ell}\), and it converges pointwise in \(W\), as \(h\to0\),
to \(\partial_\ell\partial_y^\alpha W(y,0)\).  Dominated convergence on
\((K,\mu)\) therefore gives
\[
   \partial_\ell
   \int_K \partial_y^\alpha W(y,0)\,d\mu(W)
   =
   \int_K\partial_\ell\partial_y^\alpha W(y,0)\,d\mu(W).
\]
Starting from \(\alpha=0\) and iterating this coordinate argument over the
components of a multiindex yields, for \(y\in B_R\),
\[
   \partial_y^\alpha\overline W(y)
   =
   \int_K\partial_y^\alpha W(y,0)\,d\mu(W),
\]
with the case \(\alpha=0\) understood as the definition of
\(\overline W\).  The last display also proves continuity of
\(\partial_y^\alpha\overline W\): if \(y_n\to y\) in \(B_R\), then
\(\partial_y^\alpha W(y_n,0)\to\partial_y^\alpha W(y,0)\) pointwise in
\(W\), while the integrands are dominated by \(M_{R,\alpha}\).  Hence
\(\overline W\in C^\infty_{\mathrm{loc}}\).

For the second moment, first observe that for each fixed \(y\) the map
\[
   W\mapsto W_i(y,0)W_j(y,0)
\]
is continuous on \(K\), hence Borel measurable.  It is bounded because the
scalar evaluations \(W\mapsto W_i(y,0)\) and \(W\mapsto W_j(y,0)\) have compact
images.  Hence the scalar Lebesgue integral
\[
   M_{ij}(y):=\int_K W_i(y,0)W_j(y,0)\,d\mu(W)
\]
is well defined.  The same
differentiation-under-the-integral argument applies to \(M_{ij}\).  Indeed,
on \(B_R\), Leibniz' rule writes \(\partial_y^\alpha(W_iW_j)\) as a finite sum
of products of derivatives of \(W_i\) and \(W_j\) of order at most
\(|\alpha|\).  Each factor is uniformly bounded on \(K\times B_R\) by
compactness of \(K\), and hence every product has an integrable bound
independent of \(W\).  Repeating the coordinate difference-quotient argument
used above gives
\[
   \partial_y^\alpha M_{ij}(y)
   =
   \int_K\partial_y^\alpha\bigl(W_i(y,0)W_j(y,0)\bigr)\,d\mu(W).
\]
Since \(C_{ij}=M_{ij}-\overline W_i\overline W_j\), the covariance tensor also
belongs to \(C^\infty_{\mathrm{loc}}\).

To prove that \(\overline W\) is solenoidal, take
\(\eta\in C_c^\infty(\mathbb R^3)\).  Since every \(W\in K\) is divergence
free,
\[
   \int_{\mathbb R^3}W(y,0)\cdot\nabla\eta(y)\,dy=0 .
\]
The map
\[
   (W,y)\mapsto W(y,0)\cdot\nabla\eta(y)
\]
is continuous on \(K\times\operatorname{supp}\eta\) and vanishes outside the
compact support of \(\eta\) in the \(y\)-variable.  It is therefore Borel and
absolutely bounded by
\[
   \sup_{W\in K}\|W(\cdot,0)\|_{L^\infty(\operatorname{supp}\eta)}
   \|\nabla\eta\|_{L^\infty}\mathbf 1_{\operatorname{supp}\eta}(y),
\]
which is integrable with respect to \(\mu\otimes dy\).  Fubini's theorem gives
\[
   \int_{\mathbb R^3}\overline W(y)\cdot\nabla\eta(y)\,dy
   =
   \int_K
   \int_{\mathbb R^3}W(y,0)\cdot\nabla\eta(y)\,dy\,d\mu(W)
   =
   0 .
\]
Hence \(\operatorname{div}\overline W=0\) in distributions, and therefore
classically because \(\overline W\) is smooth.

Now apply \cref{thm:statistical-identity}.  The linear term may be passed
through the \(K\)-integral by Fubini--Tonelli, since its absolute value is
bounded by
\[
   |\operatorname{supp}\phi|\,
   \sup_{W\in K}\|W(\cdot,0)\|_{L^\infty(\operatorname{supp}\phi)}
   \left\|\Delta\phi+\phi+\frac12y\cdot\nabla\phi\right\|_{L^\infty},
\]
which is finite:
\[
   \int_K\int_{\mathbb R^3}
   W\cdot\left(\Delta\phi+\phi+\frac12y\cdot\nabla\phi\right)\,dy\,d\mu(W)
   =
   \int_{\mathbb R^3}
   \overline W\cdot
   \left(\Delta\phi+\phi+\frac12y\cdot\nabla\phi\right)\,dy .
\]
For the nonlinear term, local boundedness on \(\operatorname{supp}\phi\) gives
the finite bound
\[
   |\operatorname{supp}\phi|\,
   \sup_{W\in K}\|W(\cdot,0)\|_{L^\infty(\operatorname{supp}\phi)}^2
   \|\nabla\phi\|_{L^\infty},
\]
so Fubini--Tonelli applies again.  By the definition of \(C_{ij}\),
\[
   \int_K W_iW_j\,d\mu
   =
   \overline W_i\overline W_j+C_{ij}.
\]
Therefore
\[
   \int_K\int_{\mathbb R^3}
   W_iW_j\partial_j\phi_i\,dy\,d\mu(W)
   =
   \int_{\mathbb R^3}
   \overline W_i\overline W_j\partial_j\phi_i\,dy
   +
   \int_{\mathbb R^3}
   C_{ij}\partial_j\phi_i\,dy .
\]
Substitution in \eqref{eq:statistical-identity} gives
\eqref{eq:barycenter-equation}.

Finally, we spell out the pressure-free weak stationary formulation used here.
Since \(\overline W\) is smooth and solenoidal, testing the stationary
self-similar equation
\[
   (\overline W\cdot\nabla)\overline W+\nabla \overline P
   =
   \Delta\overline W-\frac12(\overline W+y\cdot\nabla\overline W),
   \qquad \operatorname{div}\overline W=0 ,
\]
against a solenoidal \(\phi\in C_c^\infty(\mathbb R^3;\mathbb R^3)\) eliminates
the pressure term:
\[
   \int_{\mathbb R^3}\nabla\overline P\cdot\phi\,dy
   =
   -\int_{\mathbb R^3}\overline P\,\operatorname{div}\phi\,dy=0 .
\]
The same compactly supported integrations by parts used in
\cref{thm:statistical-identity} then give the pressure-free identity
\[
   \int_{\mathbb R^3}
   \overline W\cdot
   \left(\Delta\phi+\phi+\frac12 y\cdot\nabla\phi\right)\,dy
   +
   \int_{\mathbb R^3}
   \overline W_i\overline W_j\partial_j\phi_i\,dy
   =0
\]
for every solenoidal \(\phi\in C_c^\infty\).  Conversely, this is exactly what
is meant in this theorem by being stationary in the pressure-free weak
formulation: no pressure is asserted or needed in the equivalence.  Comparing
this identity with \eqref{eq:barycenter-equation} shows that the barycenter is
stationary in that weak sense if and only if the covariance term satisfies
\eqref{eq:zero-covariance} for every such test field.
\end{proof}

\begin{corollary}[Stationary support gives a stationary profile]\label{cor:stationary-support}
Let \(K\) be a compact invariant set of bounded mild ancient solutions.
Assume that there is \(L<\infty\) such that every \(W\in K\) satisfies
\[
   \sup_{\tau\in\mathbb R}\|W(\tau)\|_{L^3(\mathbb R^3)}\le L .
\]
Let \(\mu\) be an invariant probability measure on \(K\).  If \(\mu\) gives
positive mass to the set of stationary trajectories, then \(K\) contains a
stationary trajectory \(W(y,\tau)=w(y)\) whose profile satisfies
\(w\in L^3(\mathbb R^3)\).  If, in addition, the stationary trajectory is
nonzero, this contradicts the stationary self-similar theorem of
Ne\v{c}as--R\r{u}\v{z}i\v{c}ka--\v{S}ver\'ak.
\end{corollary}

\begin{proof}
Let
\[
   \mathcal S=\{W\in K:\Theta_tW=W\hbox{ for every }t\in\mathbb R\}.
\]
For \(W\in K\), it suffices to check equality for rational \(t\).  Indeed,
if \(\Theta_qW=W\) for every \(q\in\mathbb Q\), then for any
\(t\in\mathbb R\) choose \(q_n\in\mathbb Q\) with \(q_n\to t\).  Joint
continuity of the action, proved in \cref{lem:action-continuity}, gives
\(\Theta_{q_n}W\to\Theta_tW\) in \(K\), while \(\Theta_{q_n}W=W\) for all
\(n\).  Since the locally smooth topology is Hausdorff, \(\Theta_tW=W\).
Hence
\[
   \mathcal S
   =
   \bigcap_{t\in\mathbb Q}
   \{W\in K:\Theta_tW=W\}.
\]
Each set in the intersection is closed.  Indeed, for fixed \(t\) the map
\[
   K\ni W\mapsto (\Theta_tW,W)\in K\times K
\]
is continuous, and \(\{W\in K:\Theta_tW=W\}\) is the inverse image of the
diagonal in \(K\times K\).  The diagonal is closed because the locally smooth
topology is Hausdorff.  Thus \(\mathcal S\) is a countable intersection of
closed sets and is Borel.  If \(\mu(\mathcal S)>0\), then \(\mathcal S\) is
nonempty: otherwise it would be the empty Borel set and would have measure
zero.  Choose \(W\in\mathcal S\).

For this \(W\), the identity \(\Theta_tW=W\) for every \(t\) means
\[
   W(y,\tau+t)=W(y,\tau)
\]
for all \(y,\tau,t\).  Taking \(\tau=0\) gives
\(W(y,t)=W(y,0)\), so \(W\) is independent of \(\tau\).  Write
\(w(y)=W(y,0)\).  By \cref{lem:classical-local}, on every compact ball there
is a smooth pressure representative for \(W\), unique up to addition of a
function of time.  Since \(\partial_\tau W=0\), evaluating
\eqref{eq:centered} at \(\tau=0\) gives, on every compact ball, the
stationary self-similar equation
\[
   (w\cdot\nabla)w+\nabla p
   =
   \Delta w-\frac12(w+y\cdot\nabla w),
   \qquad \operatorname{div}w=0 .
\]
On overlaps of two balls the two pressure gradients agree, because both equal
the same smooth vector field
\[
   \Delta w-\frac12(w+y\cdot\nabla w)-(w\cdot\nabla)w .
\]
We now patch the local pressures.  Let \(B_m=B_m(0)\), \(m\ge1\), and choose
local pressures \(p_m\in C^\infty(B_m)\) at time \(\tau=0\).  On
\(B_m\cap B_{m+1}=B_m\) one has
\[
   \nabla(p_{m+1}-p_m)=0 .
\]
Since \(B_m\) is connected, \(p_{m+1}-p_m\) is constant on \(B_m\).  Replacing
\(p_{m+1}\) by \(p_{m+1}\) minus that constant, and proceeding inductively,
we may assume \(p_{m+1}=p_m\) on \(B_m\) for every \(m\).  The compatible
family then defines a smooth global pressure \(p\in C^\infty_{\mathrm{loc}}
(\mathbb R^3)\) by \(p=p_m\) on \(B_m\).  Hence
\[
   (w\cdot\nabla)w+\nabla p
   =
   \Delta w-\frac12(w+y\cdot\nabla w),
   \qquad \operatorname{div}w=0
\]
holds classically on every compact subset of \(\mathbb R^3\), and therefore
globally in \(\mathcal D'(\mathbb R^3)\).
The assumed uniform bound gives
\(\|w\|_{L^3(\mathbb R^3)}
=\|W(\cdot,0)\|_{L^3(\mathbb R^3)}\le L\).  The theorem of
Ne\v{c}as--R\r{u}\v{z}i\v{c}ka--\v{S}ver\'ak \cite{NRS1996} applies to smooth stationary
self-similar solutions in \(L^3(\mathbb R^3)\) and forces \(w=0\).  Therefore
any stationary trajectory in \(K\) satisfying the displayed \(L^3\) bound is
the zero trajectory, since \(W(y,\tau)=w(y)\) for all \(\tau\).  In particular
a nonzero stationary trajectory would
contradict that theorem.  This is the same stationary-rigidity interface
recorded in \cite[corollary ``Interface with stationary
rigidity'']{DuranPaper3}; the preceding paragraphs only identify the
stationary element inside the present invariant hull.
\end{proof}

\section{Conclusion}

Bounded ancient solutions generate compact trajectory hulls in the locally
smooth topology, and uniform \(L^3\)-tightness is stable on those hulls.
Compact invariant subsets support invariant probability measures obtained by
time averaging.  These measures satisfy an exact stationary statistical weak
form of the self-similar Navier--Stokes equations.  The obstruction to turning
the statistical identity into a single stationary solution is precisely the
quadratic covariance in \cref{thm:barycenter}.  A later rigidity argument must
show, under additional hypotheses, that compact nonzero ancient dynamics place
mass on stationary trajectories or otherwise make this covariance vanish.

\begin{thebibliography}{99}

\bibitem{Billingsley1999}
P.~Billingsley,
\emph{Convergence of Probability Measures},
second ed., Wiley Series in Probability and Statistics, Wiley, New York, 1999.

\bibitem{DuranPaper0}
G.~Cordero,
\emph{Local CKN Concentration and the Type I/Type II Dichotomy for
Navier--Stokes},
manuscript in this series,
\url{docs/source/type_1/unconditional_typeI_chain/paper0_local_concentration_entry_typeII.tex}.

\bibitem{DuranPaper1}
G.~Duran Ballester,
\emph{Type I Blow-up Limits for the Three-dimensional Navier--Stokes Equations
and the Centered Ancient Equation},
manuscript in this series,
\url{docs/source/type_1/paper1_typeI_blowup_limits_centered_ancient.tex}.

\bibitem{DuranPaper3}
G.~Duran Ballester,
\emph{Minimal Ancient Solutions for the Navier--Stokes Equations in
Self-similar Variables},
manuscript in this series,
\url{docs/source/type_1/paper3_minimal_ancient_solutions_reduction.tex}.

\bibitem{KryloffBogoliouboff1937}
N.~Kryloff and N.~Bogoliouboff,
\emph{La theorie generale de la mesure dans son application a l'etude des
systemes dynamiques de la mecanique non lineaire},
Ann. of Math. (2) \textbf{38} (1937), no.~1, 65--113.

\bibitem{NRS1996}
J.~Ne\v{c}as, M.~R\r{u}\v{z}i\v{c}ka, and V.~\v{S}ver\'ak,
\emph{On Leray's self-similar solutions of the Navier--Stokes equations},
Acta Math. \textbf{176} (1996), no.~2, 283--294.

\end{thebibliography}

\end{document}
